\section{Related Work}
\label{sec:related}

\subsection{Static Model Merging}
Static model merging represents the baseline approach for post-hoc multi-task integration. When multiple independent networks share a common pre-trained initialization, their parameters often reside in a common low-loss basin, allowing direct interpolation \cite{ainsworth2022git}. Early methods utilized simple parameter-space averaging, as seen in Federated Averaging (FedAvg) \cite{mcmahan2017communication} and Stochastic Weight Averaging (SWA) \cite{izmailov2018averaging}. Model Soups \cite{wortsman2022model} extended this by averaging models fine-tuned with different hyperparameters to boost generalization without inference-time cost. 

To mitigate parameter interference and conflict, several advanced merging algorithms have been proposed. Task Arithmetic \cite{ilharco2022editing} views fine-tuning directions as steerable vector fields, creating multi-task models via linear combinations of task vectors. Fisher-Weighted Averaging \cite{matena2022merging} scales parameter blends by their Fisher information matrix to preserve crucial task parameters. RegMean \cite{jin2022dataless} solves a closed-form quadratic optimization to align linear layers across fine-tuned networks. To eliminate redundant or destructive parameter updates, TIES-Merging \cite{yadav2023ties} prunes redundant sign-disagreeing parameters and applies sign-consensus, while DARE \cite{yu2023language} uses randomized delta parameter dropping to sparsify task vectors before merging. 

Despite their impressive performance, all static merging methods are fundamentally constrained: they produce a single set of static, unvarying parameters. Under high task conflict or distribution shift, no single compromise parameter state can perform optimally across all tasks, motivating the transition to localized dynamic model merging.

\subsection{Dynamic Routing and Mixture-of-Experts}
Dynamic routing and Mixture-of-Experts (MoE) architectures route input tokens or samples through specialized parameters on-the-fly. This concept originated from the adaptive mixture of local experts \cite{jacobs1991adaptive} and hierarchical MoE \cite{jordan1994hierarchical}. From a Bayesian perspective, Gaussian Processes have historically been explored as highly flexible non-parametric gating networks; for instance, \cite{rasmussen2001infinite} formulated infinite mixtures of Gaussian Process experts to route sample-wise functions, while the theoretical work of \cite{rasmussen2001occam} established how the Bayesian marginal likelihood automatically implements Occam's razor to regulate gating complexity. Modern implementations use sparsely-gated layers \cite{shazeer2017outrageously, fedus2022switch} to scale network capacity to billions of parameters, activating only a small subset of experts per input. Recent advances like Soft MoE \cite{puigcerver2023sparse} replace hard gating with soft, differentiable parameter blends to bypass routing collapse and training instability.

However, standard MoE approaches are inherently co-designed; the backbone representations and the parametric gating routing modules are trained \emph{jointly} from scratch or through extensive pre-training. This is entirely inapplicable to the post-hoc model merging regime, where the expert networks have already been fine-tuned independently, and we only possess a minuscule validation set (e.g., $N=64$ calibration samples) to calibrate our routing decisions.

\subsection{Low-Data Dynamic Merging and its Frontier}
To execute dynamic merging with minimal calibration data, several recent frameworks have explored localized routing layers. Layer-wise Linear Routing (L3-Linear) and classical parametric routers train small gating networks using standard cross-entropy on the calibration split. Under extreme data scarcity, these methods exhibit the \emph{Overfitting-Optimizer Paradox}, where high-dimensional parametric optimization overfits to spurious noise in the representations, causing test-time performance collapse.

To mitigate this, Task-Space Anchor Regularization (TSAR) \cite{rostova2025tsar} introduces an $L_2$ regularization penalty that anchors the routing weights to pre-computed expert representation centroids, while Task-Variance Regularization ($\mathcal{L}_{VR}$) in VR-Router \cite{vance2025vrrouter} dampens variance across sequential steps to prevent vectorization collapse. Quantum Wave Superposition Merging (QWS-Merge) \cite{chen2024qws} models parameter weights as quantum wave-packet states, but its complex parametric optimization remains sensitive to seed-dependent local minima. 

Recently, Parameter-Free Subspace Routing (PFSR) \cite{leland2024pfsr} introduced a non-parametric approach that projects latent representations onto task subspaces via cosine-similarity. While PFSR is immune to training loop overfitting, it relies entirely on heuristic similarity thresholds and lacks any formal Bayesian uncertainty quantification. Consequently, PFSR cannot identify out-of-distribution (OOD) test inputs, risking destructive expert parameter activation under extreme covariate shifts. 

Our proposed GP-DR bridges this critical gap. By employing a non-parametric Bayesian framework, GP-DR achieves the training-free stability of PFSR while providing mathematically grounded, closed-form posterior uncertainty quantification. Although this variance suffers from spatial collapse on the unit sphere under realistic noise (as discussed in Section~\ref{sec:method}), it remains a theoretically sound closed-form uncertainty metric that is entirely absent in prior routing systems. Table~\ref{tab:comparison_matrix} outlines how GP-DR uniquely satisfies all core requirements of robust dynamic model merging.

\begin{table}[h]
\caption{Systematic comparison of state-of-the-art dynamic model merging routing frameworks.}
\label{tab:comparison_matrix}
\vskip 0.15in
\begin{center}
\begin{small}
\begin{sc}
\resizebox{\columnwidth}{!}{%
\begin{tabular}{lcccc}
\toprule
Method & Param-Free & Training-Free & Uncertainty Quant. & Closed-Form \\
\midrule
L3-Linear \cite{rostova2025tsar} & \textsf{No}  & \textsf{No}  & \textsf{No}  & \textsf{No} \\
TSAR \cite{rostova2025tsar}       & \textsf{No}  & \textsf{No}  & \textsf{No}  & \textsf{No} \\
QWS-Merge \cite{chen2024qws}     & \textsf{No}  & \textsf{No}  & \textsf{No}  & \textsf{No} \\
PFSR \cite{leland2024pfsr}       & \textsf{Yes} & \textsf{Yes} & \textsf{No}  & \textsf{No} (Cosine Projection) \\
\textbf{GP-DR (Ours)}            & \textbf{\textsf{Yes}} & \textbf{\textsf{Yes}} & \textbf{\textsf{Yes}} & \textbf{\textsf{Yes}} \\
\bottomrule
\end{tabular}%
}
\end{sc}
\end{small}
\end{center}
\vskip -0.1in
\end{table}